\begin{claim}{D.2.54}{Proved}{A detecting class survives to the eighth stable differential}
Put $L=Y\langle7\rangle$ and use the reduced two-local
stable homotopy Atiyah--Hirzebruch spectral sequence
\[
 E^2_{p,q}=H_p(L;\pi_q^s\mathbb S)
       \Longrightarrow \pi_{p+q}^sL,
 \qquad d_r:E^r_{p,q}\longrightarrow E^r_{p-r,q+r-1}.
 \tag{D.2.54a}
\]
There is an integral class $\omega\in H_{15}(L;\ZZ)$ with
\[
 \langle\theta,\operatorname{red}_2\omega\rangle=1,
 \qquad \omega\in E^8_{15,0}.
 \tag{D.2.54b}
\]
For the integer $n$ in the actual equality
$h_Ml_4=n\widetilde\alpha\sigma'$ supplied by D.2.53,
it can be chosen so that
\[
 d_8\omega=[2n\bar\alpha]\in E^8_{7,7},
 \qquad E^2_{7,7}=H_7(L;\ZZ/16)=\ZZ/4\{\bar\alpha\}.
 \tag{D.2.54c}
\]
Here $\widetilde\alpha:S^7\to L$ is the lift of $\alpha$,
and $\bar\alpha$ denotes its integral bottom Hurewicz class
with the indicated coefficients.  The brackets in
\textup{(D.2.54c)} mean its image in the actual page-eight
quotient; they do not assert that this image is nonzero.

The only possibly nonzero incoming differentials that
form this quotient are
\[
 d_5:E^5_{12,3}\longrightarrow E^5_{7,7},\qquad
 d_6:E^6_{13,2}\longrightarrow E^6_{7,7},\qquad
 d_7:E^7_{14,1}\longrightarrow E^7_{7,7}.
 \tag{D.2.54d}
\]
Each source is its actual page group, with all earlier
differentials taken.  In particular $E^8_{7,7}$ is a
quotient of $\ZZ/4$, and the incoming $d_2$ is zero.

Let
\[
 H=E^8_{15,0}\subseteq H_{15}(L;\ZZ),\qquad
 H_0=\ker\bigl(\langle\theta,\operatorname{red}_2(-)\rangle:
                         H\longrightarrow\FF_2\bigr).
 \tag{D.2.54e}
\]
Then the complete existence test, including a change of
detecting cycle, is
\[
             h_Ml_4=0
       \quad\Longleftrightarrow\quad
             [2n\bar\alpha]\in d_8(H_0).
 \tag{D.2.54f}
\]
Consequently a proper quotient $E^8_{7,7}$ of $\ZZ/4$
would imply $h_Ml_4=0$.  A nonzero $h_Ml_4$ requires
$E^8_{7,7}=\ZZ/4$, odd $n$, and $d_8(H_0)=0$.
Neither the three incoming images in \textup{(D.2.54d)}
nor $d_8(H_0)$ is evaluated here, so D.2.4 remains open.
\end{claim}
\begin{proof}
\textbf{An actual source with only the final possible obstruction.}
Let $M_4=M(\ZZ/4,7)$ and
$\delta=\partial(2c_4):S^{14}\to M_4$.
The actual Moore map $\mu:M_4\to L$ of D.2.49,
together with D.2.53, gives
\[
                         \mu\delta=n\widetilde\alpha\sigma'.
                         \tag{D.2.54g}
\]
Choose compatible orientations and form
\[
 P=(M_4\vee S^7)\cup_{\iota_M\delta-n\iota_7\sigma'}e^{15}.
 \tag{D.2.54h}
\]
The maps $\mu$ and $\widetilde\alpha$ extend to an actual
map $g:P\to L$ by \textup{(D.2.54g)}.  If $w$ is the
oriented top integral cellular cycle, set $\omega=g_*w$.

The same actual cone comparison as D.2.52 applies with
the $S^{11}$ summand omitted.  More explicitly, pinching
the attaching sphere gives a cofiber map
\[
                     P\longrightarrow C(\delta)\vee C(-n\sigma')
                     \tag{D.2.54i}
\]
of degree one on each top relative cell.  The bottom
mod-four class can have a component on the extra $S^7$,
but its integral Bockstein has none there.  All cross
products of the two wedge summands vanish.  The actual
Stiefel-cone map $C(l_4)=V_{9,2}\to C(\delta)$ has top
degree one and bottom Moore degrees two and one.
The cup computation D.2.52j therefore gives
\[
 \langle a_7\smile\operatorname{red}_4b_8,
                  \operatorname{red}_4\omega\rangle=2,
 \qquad
 \langle\theta,\operatorname{red}_2\omega\rangle=1.
 \tag{D.2.54j}
\]
This is a cup calculation on actual spaces and their
actual cone maps; no stable splitting is used for it.

For the stable differential calculation, D.2.43 places
$\delta$ in the kernel of stabilization.  The opened
sphere convention gives $E^\infty\sigma'=2\sigma$
up to a harmless odd unit
\cite[introduction, p.~515]{EH5MorisugiMukai2000}.
Choose a stable nullhomotopy of the $M_4$ component of
the attachment in \textup{(D.2.54h)}.  Homotopy invariance
of the stable cofiber gives a filtered equivalence
\[
 \Sigma^\infty P\simeq
 \Sigma^\infty M_4\vee C(2n\sigma:S^{14}\to S^7),
 \tag{D.2.54k}
\]
identical on the old skeleton and of degree one on the
top quotient.  Thus it preserves the specified top edge
class and the cellular filtration relevant to its
differentials.  In the two-cell summand there is no cell
between dimensions seven and fifteen.  The top class
survives to page eight, and its $d_8$ is $2n\sigma$ on
the bottom seven-cell.  The separate Moore summand
contributes no differential on this top class.
Naturality for the actual map $g$ proves
\textup{(D.2.54b)--(D.2.54c)}.

We use the cellular exact-couple convention of
\cite[Part III, Section 7]{EH5Adams1974AHSS}:
$E^1$ is the cellular chain group with sphere-stem
coefficients and $d_1$ is its ordinary cellular
boundary.  There is therefore no additional unrecorded
$d_1$ in \textup{(D.2.54c)}.  The space is simply
connected, so the coefficient systems are constant.

\textbf{Every incoming source at the final target.}
Six-connectivity makes all outgoing differentials from
$(7,7)$ zero.  An incoming $d_r$, for $2\le r<8$,
would have source $(7+r,8-r)$.  The relevant coefficient
groups are
\[
 \pi_6^s\mathbb S=\ZZ/2,\quad
 \pi_5^s\mathbb S=\pi_4^s\mathbb S=0,\quad
 \pi_3^s\mathbb S=\ZZ/8,\quad
 \pi_2^s\mathbb S=\pi_1^s\mathbb S=\ZZ/2.
 \tag{D.2.54l}
\]
The zero fourth and fifth stems eliminate $d_3$ and
$d_4$ at this target.  The remaining possible $d_2$
has source $H_9(L;\ZZ/2)$, and it is zero for the
following exact naturality reason.

Let $j$ be the connective image-of-$J$ spectrum at two.
The defining fiber sequence
$j\to ko\to\Sigma^4ksp$ and the opened computation
\cite[Theorem 1.2 and Remark 1.3, pp.~2--3]
{EH5BrunerRognes2022} give
\[
                 \pi_6j=0,
       \qquad \pi_7^s\mathbb S\xrightarrow{\cong}\pi_7j=\ZZ/16
       \tag{D.2.54m}
\]
under the unit.  In the cited computation the latter
isomorphism is the four surviving $h_0$ multiples of
the image of $h_3$.  Applying the unit to the cellular
AHSS maps the source of $d_2:E^2_{9,6}\to E^2_{7,7}$
to zero and maps its target isomorphically to the
corresponding $j$-homology target.  The commuting square
therefore forces this $d_2$ to vanish.  The two-complete
convention of that source does not alter either finite
coefficient group in \textup{(D.2.54m)}.

Only the three differentials in
\textup{(D.2.54d)} remain.  This argument has not
identified their later-page sources with raw homology
groups, nor computed their images.  Successively
quotienting the target by those images proves the
assertion about $E^8_{7,7}$.

\textbf{Changing the detecting cycle and convergence.}
The sphere spectrum is connective.  Hence the bottom
coefficient row has no incoming differentials, and
$E^8_{15,0}$ is canonically a subgroup of integral
homology, as in \textup{(D.2.54e)}.  The ambiguity of
$\theta$ by $\operatorname{Sq}^1H^{14}$ evaluates to
zero on reductions of integral cycles.  Since
\textup{(D.2.54j)} gives one detecting member $\omega$,
all detecting classes on this page are exactly
$\omega+H_0$.

For $r\ge9$, a differential from $(15,0)$ would land
in filtration at most six and is zero by connectivity.
Thus the only remaining obstruction for a member of
$E^8_{15,0}$ is its $d_8$.  The connective AHSS has a
finite filtration in total degree fifteen, and its
edge image is precisely the image of
$\pi_{15}^sL\to H_{15}(L;\ZZ)$; this is the finite-range
convergence in the same cellular exact-couple construction
\cite[Part III, Section 7]{EH5Adams1974AHSS}.
It follows that a detecting stable Hurewicz class exists
if and only if
\[
                      d_8\omega\in d_8(H_0).
                      \tag{D.2.54n}
\]
Here a minus sign makes no difference to the subgroup
condition.  D.2.51 identifies this existence assertion
with $h_Ml_4=0$, proving \textup{(D.2.54f)}.

Every proper quotient of $\ZZ/4$ kills $2\bar\alpha$.
In that case the particular $\omega$ already has zero
$d_8$, and \textup{(D.2.54f)} implies $h_Ml_4=0$.
If instead $h_Ml_4\ne0$, the target must be $\ZZ/4$
and $n$ must be odd.  A nonzero subgroup of $\ZZ/4$
contains $2\bar\alpha$, so the same criterion then
forces $d_8(H_0)=0$.  These are necessary conditions,
not a computation selecting their truth values.
\end{proof}
